\subsection{A uniform presentation of the actual wild relative module}

\begin{statement}{E.75}{Proved}{The relative relation module and explicit wild kernel coordinates}
Use the actual split local source of E.74,
\[
 \mathcal D=R_3\rtimes j(\Gamma),\qquad
 D=R_3\rtimes j(P),\qquad
 P=\langle\tau,\tau'\rangle,\quad c\tau c=\tau',
 \tag{E75.1}
\]
and the cyclotomic character $\chi(\tau)=\chi(\tau')=7$.
Put $\Lambda=\mathbb Z_3[[P]]$, with its completed
group-ring topology. Conjugation by $c$ induces an
involution $\iota$ of $\Lambda$. Write
$U=\tau-7$ and $V=\tau'-7$. All modules in the following
presentation are \emph{left} compact $\Lambda$-modules;
their coefficient vectors are rows. There is an actual
continuous $\Gamma$-module isomorphism
\[
 \Phi:R_3^{\rm ab}\xrightarrow{\sim}
 \mathcal M_\chi
 :=\Lambda^2/\{\lambda(U,-V):\lambda\in\Lambda\},
 \qquad
 c(\lambda_1,\lambda_2)
   =(\iota(\lambda_2),\iota(\lambda_1)).
 \tag{E75.2}
\]
The displayed relation submodule is closed, and the map
$\lambda\mapsto\lambda(U,-V)$ is injective. The proof
constructs $\Phi$ from the actual completed Fox row of
the local source, including its inverse and its
$c$-equivariant change of basis.

Choose an actual element $r_*\in R_3$ whose image under
$\Phi$ is $(1,0)$, and put $r_*'=cr_*c$. These two
elements normally generate $R_3$ in $D$. For every
compact lattice $A=\mathcal A_r(W)$ of E.41, set
$T=T_r$, $J=J_r$ and $A^\pm=\ker(J\mp I)$. Evaluation
on this single, uniformly chosen source element gives
\[
 \begin{split}
 \operatorname{Hom}_{\Gamma,\rm cts}(R_3^{\rm ab},A)
   &\xrightarrow{\sim}\mathcal K_A,\qquad
       \theta\longmapsto\theta(r_*),\qquad\\
 \mathcal K_A
   &=\ker\bigl((I-J)(T-7I):A\longrightarrow A^-\bigr).
 \end{split}
 \tag{E75.3}
\]
For $a\in\mathcal K_A$ the inverse is the literal map
\[
 \theta_a(r)=\ell_a(\Phi(\bar r)),\qquad
 \ell_a(\lambda_1,\lambda_2)=\lambda_1a+\lambda_2Ja.
 \tag{E75.4}
\]
Thus the wild degree-one group, its coordinates and
their inverse are
\[
 \begin{split}
 H^1(G_{\mathbb Q_3},A)
   &\xrightarrow{\sim}
       \mathcal C_r(W)\oplus\mathcal K_A,\\
 [z]&\longmapsto
       ([z\circ j],z(r_*)),\\
 ([u],a)&\longmapsto
       [\,rj(g)\longmapsto\theta_a(r)+u(g)\,].
 \end{split}
 \tag{E75.5}
\]
All arrows are homeomorphisms for the compact topologies.
The source element $r_*$ and $\Phi$ are fixed once for
all coefficient lattices. Their construction uses the
actual local Fox row; it does not assign a numerical
pro-three word to that row.

More precisely, the local cochain groups in degrees
zero through two have the following actual finite
model, also valid after reduction modulo $3^m$:
\[
 A^+\xrightarrow{d^0}A\oplus A
       \xrightarrow{d^1}A^-,\qquad
 d^0(w)=((T-I)w,0),\quad
 d^1(u,a)=(I-J)(T-7I)a.
 \tag{E75.6}
\]
With the top-cell orientation specified below, its
degree-two comparison with the local invariant
coordinate of E.74 is
\[
 \begin{gathered}
 \frac{A^-}{(I-J)(T-7I)A}\xrightarrow{\sim}
 \frac{A}{(T-7I)A+(J+I)A},\\
 [y]\longmapsto[y/2],\qquad
 [x]\longmapsto[(I-J)x]\quad\text{for the inverse}.
 \end{gathered}
 \tag{E75.7}
\]

The retained program constructs exact saturated
kernel bases and inverse coordinate maps in degrees
one, two and three. Their ranks are $9,22,164$.
The nonzero Smith valuations of the map in (E75.3),
using a basis of $A^-$, are respectively
\[
 \{0^1\},\qquad
 \{0^8,1^{11},2^4\},\qquad
 \{0^{73},1^{63},2^{23},3^7\}.
 \tag{E75.8}
\]
These give exactly the degree-two obstruction groups
of E.74, including their free coordinates. The
finite model also verifies their coefficient
transitions and the eighteen finite-cohomology
order identities at $1\le m\le6$.

\emph{Manuscript: \S51.16, folios733v--736v,
scans172--178, formulas(371)--(379); Part E(c),(d)
and follow-up(4). This computes the previously
implicit saturated Hom lattice in E.74 by one
uniform source normalization. Individual nonlinear
obstruction values remain part of E.23.}
\end{statement}

\begin{proof}
\emph{An actual pair of relative normal generators.}
Use the local units $u_2=1+\lambda^2$ and
$u_3=1+\lambda^3$ from E.74, where
$\lambda=1-\zeta_3$. We use the identification
$\mu_3=\mathbb F_3$ given by the fixed root $\zeta_3$.
As elements of $k^\times/k^{\times3}$,
\[
 cu_2=u_2u_3^{-1},\qquad cu_3=u_3^{-1}.
 \tag{E75.9}
\]
Indeed $c\lambda=-\zeta_3^2\lambda$; the logarithms
modulo $\lambda^4$ give respectively
$\lambda^2-\lambda^3$ and $-\lambda^3$.
The logarithm identifies cubing on these units with
$3\lambda^2\mathcal O=\lambda^4\mathcal O$, up to a
unit, so these congruences prove the asserted Kummer
identities.

On characters with \emph{trivial} $\mathbb F_3$
coefficients the involution is the negative of
this Kummer action, because $c\zeta_3=\zeta_3^{-1}$.
Thus the character values on $cr c$ are
$(-a+b,b)$ when the values on $r$ are $(a,b)$.
The existence of these prescribed values is
the actual relative kernel calculation in
E.74: $H^1(D,\mathbb F_3)$ has Kummer basis
$\lambda,\zeta_3,u_2,u_3$, and the first
two are exactly the inflated basis of
$H^1(P,\mathbb F_3)$. Since $P$ is free,
the five-term sequence identifies
$H^1(R_3,\mathbb F_3)^P$ with the quotient
by those first two classes. The restrictions
of $\kappa_{u_2},\kappa_{u_3}$ are therefore
independent. Their joint evaluation map
$R_3\to\mathbb F_3^2$ is onto: a proper
image would give a nonzero linear relation
between those characters.
Choose $r\in R_3$ with values $(0,1)$ and put
$r'=cr c$; its values are $(1,1)$.
These two columns form a basis over $\mathbb F_3$.
The relative Frattini argument in E.74 proves that
$r,r'$ normally generate $R_3$ in $D$.
Together with $t=j(\tau)$ and $t'=j(\tau')$ they
give a minimal system of four generators of $D$.
Their involution is exactly
$t\leftrightarrow t'$ and $r\leftrightarrow r'$.

\emph{The minimal completed Fox resolution.}
We prove the needed resolution directly from the
local cohomology comparison of E.74. Put
$\Omega=\mathbb Z_3[[D]]$ and let $\mathfrak n$ be
its maximal ideal, generated by $3$ and the four
elements $t-1,t'-1,r-1,r'-1$. All maps below are
continuous maps of compact left modules. The
augmentation ideal has the epimorphism
\[
 \Omega^4\longrightarrow I_D,\qquad
 e_t\mapsto t-1,\quad e_{t'}\mapsto t'-1,\quad
 e_r\mapsto r-1,\quad e_{r'}\mapsto r'-1.
 \tag{E75.10}
\]
It is onto because the identity
$gh-1=(g-1)+g(h-1)$ expresses every group word
in its image, and that image is compact. Let $Z$
be its kernel.

Here compact Nakayama has an elementary proof.
For any compact module $M$, a set spanning
$M/\mathfrak nM$ topologically generates $M$:
otherwise a nonzero finite quotient of the remaining
compact quotient would equal its maximal-ideal
multiple. The maximal ideal of the finite acting
three-group ring is nilpotent, a contradiction.
The products in this argument are closed, since
the maximal ideal has the displayed finite
right-generating set. This proof also shows that
$\operatorname{Hom}_\Omega(M,\mathbb F_3)=0$
implies $M=0$.

Applying continuous Hom to (E75.10) and to the
augmentation sequence gives
\[
 \operatorname{Hom}_\Omega(Z,\mathbb F_3)
      =H^2(D,\mathbb F_3)=\mathbb F_3.
 \tag{E75.11}
\]
The maps on the four degree-one generators are an
isomorphism because this is a minimal system.
The bar resolution identifies the continuous
Ext groups used here with group cohomology:
evaluation on a bar basis is precisely the
continuous inhomogeneous cochain complex.
Choose a lift $e\in Z$ of a generator modulo
$\mathfrak n$. Compact Nakayama gives a
surjection $\Omega e\to Z$.
If $Z_2$ is its kernel, the same Hom sequence
and dimension shifting give
\[
 \operatorname{Hom}_\Omega(Z_2,\mathbb F_3)
   =\operatorname{Ext}^1_\Omega(Z,\mathbb F_3)
   =H^3(D,\mathbb F_3)=0.
 \tag{E75.12}
\]
The last equality is the proved local comparison
and local duality. Compact Nakayama gives $Z_2=0$.
We have proved the exact free resolution
\[
 0\longrightarrow\Omega
   \xrightarrow{d_2}\Omega^4
   \xrightarrow{d_1}\Omega
   \longrightarrow\mathbb Z_3\longrightarrow0.
 \tag{E75.13}
\]

This is the completed Fox resolution with the
chosen generators. To identify its top generator
with an actual relation, let the free pro-three
group on those four letters map to $D$, with
kernel $N$. Its complete group algebra is the
noncommutative power-series ring, as proved in
E.41. The unique last-letter decomposition gives
its free augmentation resolution. Completed
base change to $\Omega$ has first homology
$N^{\rm ab}$ and gives the exact Fox sequence
$0\to N^{\rm ab}\to\Omega^4\to I_D\to0$.
The identification sends a relation to its four
continuous Fox derivatives. It can also be
checked on the completed bar complex: choose
compatible coset representatives in the finite
quotients of the free group, and use the usual
coset decomposition in each bar factor. These
maps and their contracting homotopies are
compatible under refinement, proving the
completed Shapiro comparison. Thus $Z=N^{\rm ab}$,
and the generator $e$ has an actual relation
lift in $N$.

The involution permutes the four basis elements
of $\Omega^4$ in the two indicated pairs. On its
rank-one kernel write $c(e)=w e$.
Then $\iota(w)w=1$, where here $\iota$ denotes
the involution on $\Omega$. Modulo $\mathfrak n$
the scalar is $-1$. Indeed local duality pairs
$H^2(D,\mathbb F_3)$ with $\mu_3$, and the
field involution preserves the Brauer invariant
and inverts $\mu_3$. It therefore acts by $-1$
on this one-dimensional cohomology group and
on its dual top generator. The explicit unit
\[
 v=(1-w)/2\in\Omega^\times
 \quad\text{satisfies}\quad
 \iota(v)w=-v.
 \tag{E75.14}
\]
Replace $e$ by $ve$. Its involution is now
exactly $c(e)=-e$.

\emph{Projection to the free quotient.}
Project the four Fox coefficients to $\Lambda$.
The identity $d_1d_2=0$ says that the first two
coefficients $a_t,a_{t'}$ satisfy
$a_t(\tau-1)+a_{t'}(\tau'-1)=0$.
The augmentation ideal of the free power-series
ring is a free \emph{left} module on these two
elements, by the unique last letter. Thus they
are both zero. The exact top involution gives
the projected row
\[
 \bar d_2(1)=(0,0,a,-\iota(a)).
 \tag{E75.15}
\]
Completed base change of (E75.13) along
$\Omega\to\Lambda$ computes the homology of
$R_3$. Here the section $j$ gives the actual
coset decomposition $D=R_3j(P)$, so the
completed Shapiro maps just described use this
specified section. Since the first two
coordinates form the free augmentation
resolution of $P$, we obtain the actual map
\[
 R_3^{\rm ab}
   =\Lambda^2/\{\lambda(a,-\iota(a)):\lambda\in\Lambda\}.
 \tag{E75.16}
\]
It sends the two coordinate classes to the
images of the actual elements $r,r'$.

\emph{Finite regular coefficients and the right ideal.}
The coefficients $a,-\iota(a)$ generate the
following \emph{right} ideal:
\[
 a\Lambda+\iota(a)\Lambda=I_\chi,
 \qquad
 I_\chi=\ker\bigl(\epsilon_\chi:\Lambda\to\mathbb Z_3,
                         \ p\mapsto\chi(p)\bigr).
 \tag{E75.17}
\]
We prove equality, including its topology and
orientation. Choose a cofinal sequence of pairs
$(m,H)$, with $H\triangleleft P$ open and
$\chi(H)\equiv1\pmod{3^m}$, and put
$\Lambda_{m,H}=(\mathbb Z/3^m)[P/H]$.
Use left regular $P$ action and commuting
right multiplication. Finite local duality
\cite[I, Theorem2.1 and Corollary2.3]{E_MilneADT}
gives a natural right-module isomorphism
\[
 H^2(D,\Lambda_{m,H})
   =\bigl(\Lambda_{m,H}(-1)\bigr)_P
   \xrightarrow{\sim}\mathbb Z/3^m,
 \quad
 \sum b_p[p]\longmapsto\sum b_p\chi(p).
 \tag{E75.18}
\]
The inverse sends $1$ to $[1]$ in the
coinvariants. The right action of $p$ on the
target is multiplication by $\chi(p)$.
The transitions are the ordinary reductions
on that target. The finite degree-one
cohomology groups are finite, so their inverse
system is Mittag--Leffler. Consequently the
inverse limit of (E75.18) is the continuous
cohomology with coefficient $\Lambda$.
The finite free resolution (E75.13) also
commutes with this inverse limit and gives
\[
 \Lambda/(a\Lambda+\iota(a)\Lambda)
       =H^2(D,\Lambda)\simeq\mathbb Z_3(\chi)
       \quad\text{as right }\Lambda\text{-modules}.
 \tag{E75.19}
\]
The image of $1$ is a unit of $\mathbb Z_3$,
since it generates the right-hand module.
The annihilator of that element is exactly
$I_\chi$. This proves (E75.17); the ideal is
closed because it is a finite sum of compact
images. Rescale the top generator $e$ by the
same scalar unit so that the local invariant
in (E75.19) sends $1$ exactly to $1$.
This scalar rescaling preserves $c(e)=-e$.
It is the top-cell orientation used below.

\emph{The actual equivariant change of basis.}
The continuous substitution $p\mapsto\chi(p)^{-1}p$
is an automorphism of the completed free ring.
In power-series variables it sends
$X$ to $(X-6)/7$ and $Y$ to $(Y-6)/7$;
its inverse sends them to $6+7X,6+7Y$.
Both substitutions preserve the maximal-ideal
topology. The ordinary augmentation ideal is
right-free on $X,Y$, by unique first letter.
It follows that $I_\chi$ is right-free on
$U,V$. The unique twisted first-letter
decomposition defines a matrix $M$ by
\[
 (a,-\iota(a))=(U,V)M.
 \tag{E75.20}
\]
Equality (E75.17) says that this endomorphism
of a free right module of rank two is onto.
Its reduction modulo the maximal ideal is
therefore invertible. Lifting that inverse and
using the convergent matrix geometric series
proves $M\in\operatorname{GL}_2(\Lambda)$.
Put $S=M^{-1}\operatorname{diag}(1,-1)$.
Then
\[
 (a,-\iota(a))S=(U,-V),\qquad
 S_0\iota(S)S_0=S,\quad
 S_0=\begin{pmatrix}0&1\\1&0\end{pmatrix}.
 \tag{E75.21}
\]
For the second identity apply the involution
and swap to the first identity. Both displayed
rows change sign; uniqueness follows because
the entries of $(a,-\iota(a))$ are a right
basis of $I_\chi$.

Right multiplication of row coordinates by
$S$ gives $\Phi$ in (E75.2); its inverse is
right multiplication by $S^{-1}$.
Equation (E75.21) proves the exact involution
comparison. The relation image is compact.
It is injective, since $\lambda U=0$ forces
$\lambda=0$: the lowest nonzero word-length
coefficient of a nonzero power series is
multiplied by the nonzero constant $-6$ of
$U=X-6$. This also proves the exact sequence
$0\to\Lambda\to\Lambda^2\to R_3^{\rm ab}\to0$.

Lift $\Phi^{-1}(1,0)$ to an actual $r_*\in R_3$
and set $r_*'=cr_*c$. They form a basis of the
relative Frattini quotient, because $S$ is
invertible and both $U,V$ lie in the maximal
ideal. The relative Frattini proof again gives
normal generation. This constructs the source
pair in the statement. Its lift is a genuine
choice in $R_3$, shared by all coefficients.
For instance, to evaluate it in the original
relative basis one uses the first row of
$S^{-1}$; conversely the first row of $S$
expresses the image of $r$ in the new module.
The numerical Fox coefficients have not been
replaced by a guessed local presentation.

\emph{Cochains, inverse maps and the index-two factor.}
For a $\Gamma$-equivariant homomorphism, the
values on the two new generators are $a,Ja$.
Their one relation gives
\[
 (T-7I)a-(JTJ-7I)Ja
      =(I-J)(T-7I)a=0.
 \tag{E75.22}
\]
This proves (E75.3)--(E75.4), in both directions.
The map is saturated since the target $A^-$
is torsion-free. The cocycle restriction and
reconstruction of E.74 now give (E75.5).

One can obtain all three cochain degrees
directly from (E75.13). Lift the matrix $S$
back to $\Omega$ using the actual ring section
induced by $j$. Use the coordinate change
$(x_1,x_2)\mapsto(x_1,x_2)S$ on the last
two coordinates of $\Omega^4$; its new basis
vectors have the rows of $S^{-1}$ in the old
basis. This is invertible and $c$-equivariant.
For every coefficient module
on which $R_3$ acts trivially, the projected
Fox row is now $(0,0,U,-V)$.
The degree-zero fixed cochains are $A^+$.
The paired degree-one generators give $A^2$;
the sign on the top generator gives $A^-$.
Taking the exact $C_2$-invariant functor gives
exactly (E75.6). The same argument works for
each finite module $A/3^mA$ and its reduction
maps. The top-cell scalar has been fixed by
the $k$-local invariant in (E75.19).

Restriction from $\mathbb Q_3$ to $k$ has
degree two. Under local invariant duality it
is dual to corestriction on degree-zero Tate
duals. Thus on the coinvariants of $A(-1)$
it is the actual norm map $I-J$.
On the invariant cochain complex it sends
$y\in A^-$ to the same $y$ in the $D$
complex. Therefore its $\mathbb Q_3$ invariant
coordinate is $[y/2]$. This proves (E75.7)
with the stated sign and scalar. Directly,
$(I-J)(J+I)=0$, and
\[
 \frac{(I-J)x}{2}-x=-\frac{(I+J)x}{2};
 \qquad (I-J)(y/2)=y\quad(y\in A^-).
 \tag{E75.23}
\]
The relation maps in the two cokernels also
agree by these identities.

\emph{Exact numerical coordinates and topology.}
The program \texttt{a\_wild\_kernel.py} selects
columns of $I+J$ and $I-J$ whose reductions
form the two eigenspace bases modulo three.
Their union gives an invertible matrix over
$\mathbb Z_3$; its exact rational inverse
has denominators prime to three. In those
coordinates the target of (E75.3) is a
specified copy of $\mathbb Z_3^{\operatorname{rank}A^-}$.

Exact elimination always chooses an entry
of least three-adic valuation in the
remaining rectangle. Every subsequent
elimination factor is consequently
three-integral. The program retains the
domain transformation $Q$, its exact inverse,
the codomain transformation and the entire
diagonal identity over $\mathbb Q\cap\mathbb Z_3$.
If the rank is $s$, the last $n-s$ columns
of $Q$ give the saturated kernel basis,
and the last $n-s$ rows of $Q^{-1}$ give
its inverse coordinate map. Their product
is the identity and the original matrix
annihilates those columns exactly.
These are full integral coordinates, not
finite-precision kernel guesses.

The same computation on $(T-I)A^+$ retains
the degree-zero map. All index-two identities
in (E75.23) and both original coefficient
maps are checked exactly. The eighteen
finite-level checks verify the actual
identity
\[
 \#H^1(G_{\mathbb Q_3},A/3^mA)
  =\#(H^1(G_{\mathbb Q_3},A)/3^m)
       \cdot\#H^2(G_{\mathbb Q_3},A)[3^m]
 \quad(1\le m\le6)
 \tag{E75.24}
\]
from the two independent differentials.
The resulting kernel ranks and invariant
factors are (E75.8), matching the previously
computed groups through the literal map
(E75.7). All these modules are compact and
all basis maps are continuous, so their
continuous inverse maps also prove the
topological assertions.

Finally, for any of the actual nonlinear
geometric groups $W$, restriction to the
split source has the exact form
\[
 z(rj(g))=\vartheta(r)u(g),\quad
 \vartheta(j(g)rj(g)^{-1})
   =u(g)\,g(\vartheta(r))\,u(g)^{-1},
 \tag{E75.25}
\]
where $u$ is a $\Gamma$-cocycle and
$\vartheta:R_3\to W$ a continuous
homomorphism. Multiplication in (E75.1)
proves necessity and sufficiency. Geometric
conjugation acts on $u$ by its cocycle
formula and on $\vartheta$ by the same
inner conjugation. The module computation
above is its exact additive specialization
through $R_3\to R_3^{\rm ab}$.
For the actual geometric quotient maps
$q_r:A\to A'$, (E75.3) sends $a$ to $q_ra$;
every displayed cocycle and cochain map
commutes with this operation. No nonlinear
relation value is inferred from its
abelian image.
\end{proof}
